\section{Experimental Setup}
\label{sec:setup}

All experiments use the rosbag datasets collected in a Vicon \ac{mocap}
lab.  \autoref{tab:datasets} summarizes the three bags used for
evaluation.

\begin{table}[t]
\centering
\caption{Dataset Summary}
\label{tab:datasets}
\small
\begin{tabular}{@{}lcccc@{}}
\toprule
Bag & Dur.\ (s) & $v_{\max}$ (m/s) & $|\omega|_{\max}$ (rad/s) & Frames \\
\midrule
Slow racing & 25.8 & 3.65 & 2.62 & 258 \\
Fast racing & 18.0 & 5.66 & 9.14 & 179 \\
Backflips\textsuperscript{*} & 18.1 & $\sim$4.3 (95\%) & $\sim$10.9 (95\%) & 181 \\
\bottomrule
\multicolumn{5}{@{}p{\columnwidth}@{}}{\scriptsize \textsuperscript{*}95th
percentile; \ac{mocap} tracking artifacts dominate the instantaneous peaks.
The ${\sim}10$\,rad/s headline is this 95th-percentile sustained rate, not an
instantaneous maximum.}
\end{tabular}
\end{table}

All three regimes run the \emph{same} sliding-window solver, weighting law,
and front-end; \autoref{tab:config} collects the handful of per-regime knobs
that differ, so that every reported number can be reproduced from one place.

\begin{table}[t]
\centering
\caption{Per-Regime Sliding-Window Configuration}
\label{tab:config}
\small
\begin{tabular}{@{}lccc@{}}
\toprule
Parameter & Slow & Fast & Backflips \\
\midrule
Position knot $\Delta t_p$ (ms)        & 5  & 40  & 10 \\
Orientation knot $\Delta t_o$ (ms)     & 8  & 16  & 8 \\
LM iteration cap                       & 12 & 400 & 400 \\
Heading prior $\lambda_\psi$           & 10 & 0.6 & 10 \\
Doppler elevation corr.\ $b$ (m/s)     & --- & --- & $-1.5$ \\
Position-init anchor $\lambda_{p_0}$   & --- & --- & 0.5 \\
\bottomrule
\multicolumn{4}{@{}p{0.92\columnwidth}@{}}{\footnotesize Shared across all
regimes: 3\,s window, 0.3\,s stride, unscaled Markov-blanket marginalization
prior, the single dynamics-adaptive weighting law (\autoref{sec:models}),
RANSAC ego-velocity prefilter (\SI{0.15}{\metre\per\second} inlier),
sensor-only initialization, and frozen extrinsic pitch \SI{27.5}{\degree}.
``---'' marks the shared default: uncapped-by-data iterations, and the
flight-regime elevation correction / position tether off.}
\end{tabular}
\end{table}

\textbf{\ac{imu} rate:} $\sim$993\,Hz (full rate, no downsampling for C++ solver).
\textbf{Evaluation:} Start-anchored rigid alignment to \ac{mocap} ground
truth: the alignment rotation $\mathbf{R}_\mathrm{align} =
\mathbf{R}_{\mathrm{gt},0}\,\mathbf{R}_{\mathrm{est},0}^\top$
and translation $\mathbf{t}_\mathrm{align} =
\mathbf{p}_{\mathrm{gt},0} - \mathbf{R}_\mathrm{align}\,
\mathbf{p}_{\mathrm{est},0}$ are fixed at frame~0 (no scale).
The final 3\,s is trimmed.  Velocity \ac{rmse} is computed from the
analytical first derivative of the position B-spline, not from finite
differences of the estimated trajectory.
\emph{Position drift} (\%) is position \ac{rmse} normalised by the total
GT path length, the standard odometry metric for a
dead-reckoning system with no loop closure.
Start-anchored alignment is harsher than the whole-trajectory
SE(3)/Umeyama fit (the EVO/KITTI default) used by most odometry papers,
because it cannot retro-fit a constant transform to cancel accumulated
drift: it fixes the frame once, at the start, exactly as a deployed
estimator must.  We adopt it precisely for this causal,
deployment-relevant reading, and apply it identically to every system
compared.  Where we compare against externally published figures
(\autoref{sec:baselines}) we additionally report whole-trajectory-aligned
numbers so the baselines can be cross-checked against their own papers.
\emph{Translational \ac{rpe}} (relative pose error, \autoref{fig:rpe}) is
the KITTI-standard per-segment displacement error averaged over all
sub-segments of a given accumulated GT length.
Because the start-anchored alignment offsets cancel within every
relative displacement, \ac{rpe} is alignment-independent.
For sliding window, \emph{live (leading-edge)} \ac{rmse} is computed by
snapshotting the estimate at $[t_\text{end}{-}\text{stride},\, t_\text{end}]$
of each window at 50\,Hz before any subsequent refinement.  The
\emph{settled} \ac{rmse} evaluates the same trajectory after all windows
have refined each point.

\textbf{Solver implementation.}
The Ceres LM problem is solved with \texttt{SPARSE\_NORMAL\_CHOLESKY}
(SuiteSparse), up to 400 iterations.
All three factor types (gyroscope, accelerometer, and radar Doppler)
use hand-derived analytic Jacobians based on the basalt
cumulative-spline Jacobian struct; Ceres automatic differentiation was
eliminated to avoid Jet-arithmetic overhead over the large set of
spline control-point parameters per factor.
In the sliding-window solver the linear solve and the
marginalization-prior recompute (\texttt{compute\_prior}, evaluated
\emph{every} LM iteration, as distinct from the ${\sim}10$\,ms one-time
prior \emph{formation} of \autoref{sec:sw}) are the dominant costs;
Jacobian evaluation accounts for roughly 25--30\% of per-window wall time.
